% Auto-generated from meeting_2566-10-20_curriculum-draft-2.md (กบ.วช. format v2)
\documentclass{meetingminutes}
\meetingcommittee{คณะกรรมการบริหารหลักสูตร วิศวกรรมคอมพิวเตอร์และปัญญาประดิษฐ์ (บัณฑิตศึกษา)}
\meetingnumber{๒}{๒๕๖๖}
\meetingdate{วันที่ ๒๐ ตุลาคม พ.ศ. ๒๕๖๖}
\meetingplace{ห้องประชุมคณะเทคโนโลยีอุตสาหกรรม}
\meetingstart{๑๓.๐๐ น.}
\meetingend{๑๕.๕๐ น.}

\begin{document}
\maketitle

\psection{ผู้มาประชุม}
\begin{participants}
  \pmember{1}{รองศาสตราจารย์ ดร.อิสระ อินจันทร์}{ที่ปรึกษาหลักสูตร}{ประธาน}
  \pmember{2}{รองศาสตราจารย์ ดร.กันต์ อินทุวงศ์}{ที่ปรึกษาหลักสูตร}{รองประธาน}
  \pmember{3}{ผู้ช่วยศาสตราจารย์ ดร.พิศิษฐ์ นาคใจ}{อาจารย์ประจำหลักสูตร}{กรรมการ}
  \pmember{4}{ผู้ช่วยศาสตราจารย์ ดร.กาญจนา ดาวเด่น}{อาจารย์ประจำหลักสูตร (เลขานุการ)}{กรรมการ/เลขานุการ}
  \pmember{5}{ผู้ช่วยศาสตราจารย์ ดร.พัชรี มณีรัตน์}{อาจารย์ประจำหลักสูตร}{กรรมการ}
  \pmember{6}{ผู้ช่วยศาสตราจารย์ ดร.พิทักษ์ คล้ายชม}{รองคณบดีฝ่ายวิชาการ คณะเทคโนโลยีอุตสาหกรรม}{กรรมการ}
  \pmember{7}{ผู้ช่วยศาสตราจารย์ ดร.โสกณ วิริยะรัตนกุล}{อาจารย์ประจำหลักสูตร}{กรรมการ}
  \pmember{8}{ผู้ช่วยศาสตราจารย์ ดร.ธัชชัย อยู่ยิ่ง}{อาจารย์ประจำหลักสูตร}{กรรมการ}
  \pmember{9}{ผู้ช่วยศาสตราจารย์ ดร.อภิศักดิ์ พรหมฝ่าย}{อาจารย์ประจำหลักสูตร}{กรรมการ}
\end{participants}

\psection{ผู้ไม่มาประชุม}
\noindent ไม่มี\par

\startmeeting
\openingchair{รองศาสตราจารย์ ดร.อิสระ อินจันทร์}{กล่าวเปิดการประชุมและดำเนินการประชุมตามระเบียบวาระ ดังนี้}

\agenda{๑}{เรื่องแจ้งให้ที่ประชุมทราบ}
ประธานแจ้งความคืบหน้าการพัฒนาหลักสูตรหลังการประชุมครั้งที่ 1/2566 โดยคณะทำงานได้ศึกษาหลักสูตรอ้างอิงจากมหาวิทยาลัยชั้นนำ ได้แก่ จุฬาลงกรณ์มหาวิทยาลัย มหาวิทยาลัยเทคโนโลยีพระจอมเกล้าธนบุรี และมหาวิทยาลัยเชียงใหม่ รวมถึงศึกษามาตรฐานคุณวุฒิระดับบัณฑิตศึกษา กมอ. พ.ศ. 2565 และได้จัดทำ กรอบ PLO เบื้องต้นจากการวิเคราะห์เอกสาร (document analysis) เพื่อใช้เป็นฐานในการรับฟังความคิดเห็นจากผู้มีส่วนได้ส่วนเสียต่อไป

ประธานย้ำว่ากรอบ PLO ที่จะนำเสนอในการประชุมครั้งนี้เป็นเพียง ข้อเสนอเบื้องต้น (preliminary framework) ซึ่งยังไม่ใช่ PLO ฉบับสมบูรณ์ และจะต้องผ่านกระบวนการรับฟังความคิดเห็น stakeholder ในเดือนพฤศจิกายน 2566 ก่อนจึงจะสามารถกำหนด PLO อย่างเป็นทางการได้ ตามข้อกำหนด AUN-QA R1.4


\agenda{๒}{รับรองรายงานการประชุมครั้งที่ 1/2566}
ที่ประชุมรับรองรายงานการประชุมครั้งที่ 1/2566 (15 กันยายน 2566) โดยไม่มีการแก้ไข


\agenda{๓}{เรื่องสืบเนื่อง}
\subagenda{๓.๑}{รายงานกรอบ PLO เบื้องต้นจากการวิเคราะห์เอกสาร (Pre-Stakeholder Draft)}
ผศ.ดร.พิศิษฐ์ นำเสนอ กรอบ PLO เบื้องต้น จำนวน 5 ข้อ ซึ่งได้จากการวิเคราะห์: (1) มาตรฐานคุณวุฒิ กมอ. (2) หลักสูตรอ้างอิงจากสถาบันชั้นนำ และ (3) แนวโน้มอุตสาหกรรม AI และวิศวกรรมคอมพิวเตอร์ในประเทศ โดยยังไม่ได้ผ่านการรับฟัง stakeholder พร้อมระดับ Bloom's Taxonomy ดังนี้:

\begin{longtable}{|p{2cm}|p{7cm}|p{2.5cm}|p{2cm}|}
  \hline
  \textbf{PLO} & \textbf{ร่างผลลัพธ์การเรียนรู้} & \textbf{Bloom's Level} & \textbf{ประเภท} \\
  \hline
  PLO1 & ใช้ซอฟต์แวร์หรือฮาร์ดแวร์ทางคอมพิวเตอร์ในการสร้างสรรค์ผลงาน องค์ความรู้ และนวัตกรรม & Create (6) & SSLO \\
  \hline
  PLO2 & ประยุกต์ใช้ความรู้ในสาขาวิชาวิศวกรรมคอมพิวเตอร์เพื่อแก้ไขปัญหาความต้องการขององค์กร ชุมชน และท้องถิ่น & Apply (3) & SSLO \\
  \hline
  PLO3 & สร้างองค์ความรู้ทางด้านวิศวกรรมคอมพิวเตอร์และปัญญาประดิษฐ์ที่ตอบสนองความต้องการขององค์กร & Create (6) & SSLO \\
  \hline
  PLO4 & ปฏิบัติตามหลักจรรยาบรรณวิชาชีพวิศวกรรมคอมพิวเตอร์ & Evaluate (5) & GLO \\
  \hline
  PLO5 & สามารถคิดวิเคราะห์ การสื่อสาร การทำงานร่วมกัน การปรับตัว และการเรียนรู้ตลอดชีวิต & Analyze (4) & GLO \\
  \hline
\end{longtable}

ที่ประชุมอภิปราย PLOs ร่างและเห็นว่า PLO1-PLO3 เป็น Subject-Specific Learning Outcomes (SSLO) ส่วน PLO4-PLO5 เป็น Generic Learning Outcomes (GLO) สอดคล้องกับ R1.3 ของ AUN-QA

\resolution{รับทราบกรอบ PLO เบื้องต้นทั้ง 5 ข้อในฐานะสมมติฐานเพื่อใช้ในการรับฟัง stakeholder เท่านั้น ยังไม่ถือว่าอนุมัติ PLO อย่างเป็นทางการ มอบหมายให้นำกรอบนี้ไปเป็นแนวทางถามความเห็นในการประชุม stakeholder consultation เดือนพฤศจิกายน 2566 และปรับปรุง PLO ให้สอดคล้องกับ stakeholder needs ก่อนนำกลับมาพิจารณาอนุมัติ}

\agenda{๔}{เรื่องเสนอเพื่อพิจารณา}
\subagenda{๔.๑}{โครงสร้างรายวิชาและหน่วยกิตในหลักสูตร}
ที่ประชุมพิจารณาโครงสร้างรายวิชา โดยกำหนดหมวดวิชาดังนี้: - วิชาเฉพาะด้านบังคับ (6 วิชา, 18 หน่วยกิต) - วิชาเฉพาะด้านเลือก (เลือกตามแผน) - วิทยานิพนธ์/สารนิพนธ์ (ตามแผน) - วิชาเสริม — ภาษาอังกฤษและการเขียนวิทยานิพนธ์ (6 หน่วยกิต ไม่นับรวม)

\resolution{เห็นชอบโครงสร้างหน่วยกิตตามที่เสนอ มอบหมายให้คณะทำงานยกร่างรายละเอียดวิชาทั้งหมด (course description + CLO) และตาราง Curriculum Mapping}
\subagenda{๔.๒}{การเชื่อมโยง CLO กับ PLO (Curriculum Mapping)}
ที่ประชุมอภิปรายหลักการ Constructive Alignment และ OBE (Outcome-Based Education) โดยย้ำว่าทุก CLO ของทุกรายวิชาต้องเชื่อมโยงกับ PLO อย่างน้อย 1 ข้อ และการออกแบบ CLO ต้องใช้ Action Verb ที่เหมาะสมกับระดับ Bloom's Taxonomy

\resolution{มอบหมายให้อาจารย์แต่ละท่านออกแบบ CLO ของวิชาที่รับผิดชอบ โดยนำเสนอในการประชุมครั้งถัดไป}

\agenda{๕}{เรื่องอื่น ๆ}
ที่ประชุมเห็นชอบให้มีการจัดประชุมรับฟังความคิดเห็นจากผู้มีส่วนได้ส่วนเสียภายนอก ในเดือนพฤศจิกายน 2566


\adjournmeeting

\begin{sigblock}
\typist{ผู้ช่วยศาสตราจารย์ ดร.กาญจนา ดาวเด่น}
\reviewer{รองศาสตราจารย์ ดร.อิสระ อินจันทร์}{กรรมการ}
\chairsig{รองศาสตราจารย์ ดร.อิสระ อินจันทร์}{ประธานการประชุม}
\end{sigblock}

\end{document}
